\documentclass[a4paper, 12pt]{article}

\usepackage[utf8]{inputenc}
\usepackage{enumitem}
\usepackage{amsmath}
\usepackage{amssymb}

\title{E.B.1.1}
\author{Lorenzo Burello}
\date{\today}

\begin{document}
\maketitle

\section{Definizione del problema}
\subsection{Consegna}
Modellare, mediante una formula $\phi$ in logica del primo ordine,
la seguente affermazione, scegliendo opportuni insiemi di simboli
di predicato e/o funzione: \newline
\begin{quote}
    Nessuno studente supererà un qualunque esame se è ansioso o non ne studia il programma.
\end{quote}
\subsection{Definizione dell'alfabeto}

$\mathcal{P}$:
\begin{itemize}[label=-]
    \item Studente/1
    \item Esame/1
    \item Ansioso/1
    \item Programma\_studiato/2
    \item Esame\_superato/2
\end{itemize}

\subsection{Riformulazione della domanda}
\begin{quote}
    \textbf{Nessuno} studente supererà un qualunque esame se è ansioso o non ne studia il programma.
\end{quote}
Può diventare
\begin{quote}
    \textbf{Ogni} studente \textbf{non} supererà un qualunque esame se è ansioso o non ne studia il programma.
\end{quote}

\section{Creazione di $\phi$}

\[
\phi =
\begin{array}{c}
Studente(s) \\
\land \\
Esame(e) \\
\land \\
(Ansioso(s) \lor \lnot Programma\_studiato(s,e))
\end{array}
\rightarrow
\lnot Esame\_superato(s,e)
\]

\section{Determinazione interpretazioni}
\subsection{$ M \models \phi $}

Sia $D = \{\alpha, \beta, \gamma \}$ il dominio di $M$:
\begin{itemize}[label=-]
    \item $ M(Studente) = \{ \alpha, \beta \} $
    \item $ M(Ansioso) = \{ \beta \} $
    \item $ M(Esame) = \{ \gamma \} $
    \item $ M(Programma\_studiato)=\{ (\alpha, \gamma),( \beta, \gamma) \} $
    \item $ M(Esame\_superato) = \{ (\alpha, \gamma) \}$
\end{itemize}

\subsection{$ I \not \models \phi$}
Sia $D = \{\alpha, \beta, \gamma \}$ il dominio di $I$:
\begin{itemize}[label=-]
    \item $ I(Studente) = \{ \alpha, \beta \} $
    \item $ I(Ansioso) = \{ \beta \} $
    \item $ I(Esame) = \{ \gamma \} $
    \item $ I(Programma\_studiato)=\{ (\alpha, \gamma),( \beta, \gamma) \} $
    \item $ I(Esame\_superato) = \{ (\beta, \gamma) \}$
\end{itemize}
\end{document}